% TOP_PROBLEM: {"id": 30006987, "title": "Nagata Conjecture on Plane Curve Linear Systems", "category_id": 5, "category": {"id": 5, "name": "algebraic_geometry", "display_name": "Algebraic Geometry", "description": "Geometric objects defined by polynomial equations.", "slug": "algebraic-geometry", "order_index": 5, "created_at": "2026-07-31T15:26:25.671Z"}, "status": "open"}
% FIRST_POSTED: 2026-10-06
% ATTEMPT_STATUS: solved
% OPENAI_CLAIM: {"schema_version":1,"source_repo":"https://github.com/openai/math","source_commit":"adc7f1241b42e322a6451854ab7e4b4c146bf78a","release_date":"2026-10-06","match":"full","resolution":"proved","independently_reviewed":false,"second_pass":"agreed","families":[{"family":"039","title":"Nagata’s conjecture and maximal Seshadri constants","subject":"Algebraic and complex geometry","manuscripts":[{"title":"Maximal Multipoint Seshadri Constants in Higher Dimensions","pdf_path":"preprints/Maximal-Multipoint-Seshadri-Constants-in-Higher-Dimensions-October-5-2026/main.pdf","readme_path":"preprints/Maximal-Multipoint-Seshadri-Constants-in-Higher-Dimensions-October-5-2026/README.md","date":"2026-10-05","theorem_ref":"thm:main (sections/introduction.tex:29)","statement_sources":[{"path":"preprints/Maximal-Multipoint-Seshadri-Constants-in-Higher-Dimensions-October-5-2026/build/sections/introduction.tex","line":29,"label":"thm:main"}]},{"title":"Maximal Seshadri constants on arbitrary polarized surfaces","pdf_path":"preprints/Maximal-Seshadri-Constants-on-Arbitrary-Polarized-Surfaces-September-23-2026/main.pdf","readme_path":"preprints/Maximal-Seshadri-Constants-on-Arbitrary-Polarized-Surfaces-September-23-2026/README.md","date":"2026-09-23","theorem_ref":"thm:main (sections/introduction.tex:31)","statement_sources":[{"path":"preprints/Maximal-Seshadri-Constants-on-Arbitrary-Polarized-Surfaces-September-23-2026/build/sections/introduction.tex","line":31,"label":"thm:main"}]},{"title":"Nagata's conjecture for plane curves","pdf_path":"preprints/Nagatas-Conjecture-for-Plane-Curves-September-23-2026/main.pdf","readme_path":"preprints/Nagatas-Conjecture-for-Plane-Curves-September-23-2026/README.md","date":"2026-09-23","theorem_ref":"thm:main (sections/introduction.tex:31)","statement_sources":[{"path":"preprints/Nagatas-Conjecture-for-Plane-Curves-September-23-2026/build/sections/introduction.tex","line":31,"label":"thm:main"}]}],"lean":{"doc_path":"lean/docs/039.md","comparators":["lean/ComparatorChallenges/MaximalSeshadriConstants.json","lean/ComparatorChallenges/MaximalSeshadriConstants.lean","lean/ComparatorChallenges/Nagata.json","lean/ComparatorChallenges/Nagata.lean"],"declarations":["OAI.MaximalSeshadri.Geometry.maximalSeshadriConstants","OAI.Nagata.nagata_conjecture"],"covers_main_theorem":false},"reasoning_trace":null}]}
\documentclass[11pt]{article}
\usepackage[margin=25mm]{geometry}
\usepackage{fontspec}
\setmainfont{Latin Modern Roman}
\usepackage{amsmath,amssymb,mathtools}
\usepackage{unicode-math}
\setmathfont{Latin Modern Math}
\usepackage{enumitem}
\usepackage{xurl,hyperref,bookmark}
\hypersetup{colorlinks=true,urlcolor=blue}
\setcounter{secnumdepth}{0}
\setlength{\parindent}{0pt}
\setlength{\parskip}{5pt}
\setlength{\emergencystretch}{3em}
\begin{document}
\section*{Nagata Conjecture on Plane Curve Linear Systems}

\subsection{OpenAI's claim}
Manuscript 1 (family 039, 2026-10-05): In every dimension at least three, the ordinary multipoint Seshadri constant is claimed to reach $(L^n/r)^{1/n}$ at very general points for every sufficiently large $r$. Manuscript 2 (family 039, 2026-09-23): For any polarized smooth complex projective surface, the equal-weight multipoint Seshadri constant is claimed to reach $\sqrt{L^2/r}$ for all integers beyond a polarization-dependent threshold. Manuscript 3 (family 039, 2026-09-23): At $r\ge10$ very general complex plane points, every effective curve is claimed to have total assigned multiplicity below its degree times $\sqrt r$, including reducible curves and unequal assignments.

\subsection{Mathematical statement of the claimed result}
\paragraph{Mathematical claim 1.} For each $r\ge10$, OpenAI claims a nonempty very-general locus of ordered distinct points $(p_1,\ldots,p_r)$ in $\mathbb P^2_{\mathbb C}$ such that every nonzero effective plane curve $C$ of degree $d\ge1$ and integers $m_i\ge0$ with $\operatorname{mult}_{p_i}C\ge m_i$ satisfy
\[
d\sqrt r>\sum_{i=1}^r m_i.
\]
Very general means excluding a union, indexed by a countable set, of proper algebraic closed subsets of the distinct-point configuration space. Multiplicities may be unequal, curves may be reducible, and the strict inequality starts at ten points.

\paragraph{Mathematical claim 2.} Let $X$ be smooth integral projective over $\mathbb C$, with $n=\dim X\ge3$, and $L$ ample. OpenAI claims an integer $r_0(X,L)$ such that for every $r\ge r_0$ and very general ordered distinct $p_1,\ldots,p_r$,
\[
\varepsilon(X,L;\mathbf p)=\left(\frac{L^n}{r}\right)^{1/n}.
\]
The multipoint Seshadri constant is the supremum of $t\ge0$ for which $\pi^*L-t\sum_iE_i$ is nef on the blow-up at these points. For each $r$, the exceptional tuples lie in a union, indexed by a countable set, of proper algebraic closed subsets with nonempty complement. The threshold depends on $(X,L)$ and is not specified uniformly.

\paragraph{Mathematical claim 3.} For any smooth integral projective complex surface $S$ and ample line bundle $L$, OpenAI claims $r_0(S,L)$ such that every $r\ge r_0$ and very general tuple of distinct points give, on the blow-up $\pi$ with exceptional curves $E_i$,
\[
\pi^*L-\sqrt{L^2/r}\sum_{i=1}^rE_i\text{ nef},\qquad
\varepsilon(S,L;\mathbf p)=\sqrt{L^2/r}.
\]
The constant is the supremum of $t$ making $\pi^*L-t\sum E_i$ nef. Very general excludes countably many proper closed subsets of the distinct-point configuration space; the statement is for sufficiently many points, not every prescribed small $r$.

\subsection{Outline of the argument}
Manuscript 1 (family 039, 2026-10-05): Graded exponent supports are reshaped into a volume-preserving simplex compatible with section multiplication. Full density forces every sufficiently deep interior lattice point to occur, allowing torus interpolation to separate the needed jets; very-general tuples inherit all the finite tests.

Manuscript 2 (family 039, 2026-09-23): Finite jet tests are transferred to an algebraic torus through nodal-curve degeneration. A varying exponent region is tuned to a primitive lattice direction and a subgroup with exactly $r$ points, giving injectivity for all required orders rather than only a fixed multiplicity.

Manuscript 3 (family 039, 2026-09-23): Incidence loci and cyclic products reduce a hypothetical failure to equal multiplicities at arbitrary configurations. Degeneration toward a cubic produces finite jet constraints; colliding the remaining points and taking limits of elliptic theta bases turns the contradiction into a full-rank polynomial evaluation problem.

\subsection{Scope relative to this problem}
Manuscript 1 (family 039, 2026-10-05): Eventual maximal multipoint constants do not cover every plane point-count r>=10 in Nagata. The surface companion in this build also has an input-dependent eventual threshold.

Manuscript 2 (family 039, 2026-09-23): Maximal Seshadri constants imply a plane-curve slope bound only for sufficiently many points with an unspecified threshold, and give a non-strict bound; our Nagata record requires strictness for every r at least ten.

Manuscript 3 (family 039, 2026-09-23): The main theorem gives the exact strict degree/multiplicity inequality for every r at least ten, all unequal nonnegative multiplicities and every nonzero effective plane curve at very general complex points.

\subsection{Formal verification}
For family 039, the scope page documents the following selections. The selected plane-curve inequality allows unequal multiplicities and repeated components at every allowed point count. A separate statement gives eventual maximality on each polarized complex surface, with a surface-dependent threshold; the higher-dimensional and positive-characteristic companions are not selected here. Manuscript 1 is not among the papers selected by that page, so the family link does not formalize its main conclusion. The documented statement covers the main claim of manuscript 2. The documented statement covers the main claim of manuscript 3. Its selected declaration also covers the clause used for this problem. Comparator files: MaximalSeshadriConstants.lean, Nagata.lean. Declarations: OAI.MaximalSeshadri.Geometry.maximalSeshadriConstants, OAI.Nagata.nagata\_conjecture.

This site has not executed the Lean files or checked their dependencies. These descriptions report the scope of the linked declarations, rather than a new verification of the proofs.

\subsection{Status caveat}
OpenAI attributes these results to an internal model. Its release notice cautions that results lacking full formalization can contain mistakes. This site has not independently assessed the proofs; a solved label records the scope of the reported claim.

\subsection{Sources}
\begin{itemize}
\item OpenAI, \emph{Maximal Multipoint Seshadri Constants in Higher Dimensions}, 2026-10-05; citation key \texttt{OAI:Maximal-Multipoint-Seshadri-Constants-in-Higher-Dimensions-October-5-2026}. \href{https://github.com/openai/math/blob/adc7f1241b42e322a6451854ab7e4b4c146bf78a/preprints/Maximal-Multipoint-Seshadri-Constants-in-Higher-Dimensions-October-5-2026/main.pdf}{PDF (GitHub)}; \href{https://github.com/openai/math/raw/adc7f1241b42e322a6451854ab7e4b4c146bf78a/preprints/Maximal-Multipoint-Seshadri-Constants-in-Higher-Dimensions-October-5-2026/main.pdf}{PDF (direct)}; \href{https://github.com/openai/math/blob/adc7f1241b42e322a6451854ab7e4b4c146bf78a/preprints/Maximal-Multipoint-Seshadri-Constants-in-Higher-Dimensions-October-5-2026/README.md}{Manuscript page and citation}; \href{https://github.com/openai/math/blob/main/preprints/Maximal-Multipoint-Seshadri-Constants-in-Higher-Dimensions-October-5-2026/main.pdf}{Latest version}; Statement source: \href{https://github.com/openai/math/blob/adc7f1241b42e322a6451854ab7e4b4c146bf78a/preprints/Maximal-Multipoint-Seshadri-Constants-in-Higher-Dimensions-October-5-2026/build/sections/introduction.tex\#L29}{thm:main}.
\item OpenAI, \emph{Maximal Seshadri constants on arbitrary polarized surfaces}, 2026-09-23; citation key \texttt{OAI:Maximal-Seshadri-Constants-on-Arbitrary-Polarized-Surfaces-September-23-2026}. \href{https://github.com/openai/math/blob/adc7f1241b42e322a6451854ab7e4b4c146bf78a/preprints/Maximal-Seshadri-Constants-on-Arbitrary-Polarized-Surfaces-September-23-2026/main.pdf}{PDF (GitHub)}; \href{https://github.com/openai/math/raw/adc7f1241b42e322a6451854ab7e4b4c146bf78a/preprints/Maximal-Seshadri-Constants-on-Arbitrary-Polarized-Surfaces-September-23-2026/main.pdf}{PDF (direct)}; \href{https://github.com/openai/math/blob/adc7f1241b42e322a6451854ab7e4b4c146bf78a/preprints/Maximal-Seshadri-Constants-on-Arbitrary-Polarized-Surfaces-September-23-2026/README.md}{Manuscript page and citation}; \href{https://github.com/openai/math/blob/main/preprints/Maximal-Seshadri-Constants-on-Arbitrary-Polarized-Surfaces-September-23-2026/main.pdf}{Latest version}; Statement source: \href{https://github.com/openai/math/blob/adc7f1241b42e322a6451854ab7e4b4c146bf78a/preprints/Maximal-Seshadri-Constants-on-Arbitrary-Polarized-Surfaces-September-23-2026/build/sections/introduction.tex\#L31}{thm:main}.
\item OpenAI, \emph{Nagata's conjecture for plane curves}, 2026-09-23; citation key \texttt{OAI:Nagatas-Conjecture-for-Plane-Curves-September-23-2026}. \href{https://github.com/openai/math/blob/adc7f1241b42e322a6451854ab7e4b4c146bf78a/preprints/Nagatas-Conjecture-for-Plane-Curves-September-23-2026/main.pdf}{PDF (GitHub)}; \href{https://github.com/openai/math/raw/adc7f1241b42e322a6451854ab7e4b4c146bf78a/preprints/Nagatas-Conjecture-for-Plane-Curves-September-23-2026/main.pdf}{PDF (direct)}; \href{https://github.com/openai/math/blob/adc7f1241b42e322a6451854ab7e4b4c146bf78a/preprints/Nagatas-Conjecture-for-Plane-Curves-September-23-2026/README.md}{Manuscript page and citation}; \href{https://github.com/openai/math/blob/main/preprints/Nagatas-Conjecture-for-Plane-Curves-September-23-2026/main.pdf}{Latest version}; Statement source: \href{https://github.com/openai/math/blob/adc7f1241b42e322a6451854ab7e4b4c146bf78a/preprints/Nagatas-Conjecture-for-Plane-Curves-September-23-2026/build/sections/introduction.tex\#L31}{thm:main}.
\item Family 039: \href{https://github.com/openai/math/blob/adc7f1241b42e322a6451854ab7e4b4c146bf78a/lean/docs/039.md}{Lean scope}.
\item \href{https://github.com/openai/math/blob/adc7f1241b42e322a6451854ab7e4b4c146bf78a/lean/ComparatorChallenges/MaximalSeshadriConstants.json}{Comparator MaximalSeshadriConstants.json}.
\item \href{https://github.com/openai/math/blob/adc7f1241b42e322a6451854ab7e4b4c146bf78a/lean/ComparatorChallenges/MaximalSeshadriConstants.lean}{Comparator MaximalSeshadriConstants.lean}.
\item \href{https://github.com/openai/math/blob/adc7f1241b42e322a6451854ab7e4b4c146bf78a/lean/ComparatorChallenges/Nagata.json}{Comparator Nagata.json}.
\item \href{https://github.com/openai/math/blob/adc7f1241b42e322a6451854ab7e4b4c146bf78a/lean/ComparatorChallenges/Nagata.lean}{Comparator Nagata.lean}.
\item \href{https://github.com/openai/math/blob/adc7f1241b42e322a6451854ab7e4b4c146bf78a/CONTENTS.md}{OpenAI catalogue entry}.
\end{itemize}
\end{document}
